\documentclass[a4paper,11pt,titlepage]{article}

\usepackage[utf8]{inputenc}
\usepackage[spanish]{babel}
\usepackage{fullpage}
\usepackage[table]{xcolor}
\usepackage{tabularx}
\usepackage{csquotes}
\usepackage[
    backend=biber, 
    style=numeric, 
    sorting=none, 
    maxcitenames=2, % Si hay 3 o más, pone et al.
    mincitenames=1  % Al recortar, deja solo al primer autor
]{biblatex}
\addbibresource{./bib/reference.bib}
\usepackage[colorlinks=true, allcolors=blue]{hyperref}
% Datos generales del proyecto.
% Editar este archivo al crear una tesis nueva desde la plantilla.

\newcommand{\tesisAutor}{Nombre Apellido}
\newcommand{\tesisAutorCorto}{Apellido}
\newcommand{\tesisEmail}{correo@ejemplo.com}
\newcommand{\tesisLU}{000/00}
\newcommand{\tesisTitulo}{Titulo de la tesis}
\newcommand{\tesisTituloCorto}{Titulo corto}
\newcommand{\tesisCarrera}{Licenciatura en Ciencias de Datos}
\newcommand{\tesisDirector}{Nombre de la direccion}
\newcommand{\tesisDirectorTitulo}{Dr./Dra. Nombre de la direccion}
\newcommand{\tesisDirectorEmail}{direccion@ejemplo.com}
\newcommand{\tesisCodirector}{Nombre de la codireccion}
\newcommand{\tesisCodirectorTitulo}{Dr./Dra. Nombre de la codireccion}
\newcommand{\tesisCodirectorEmail}{codireccion@ejemplo.com}
\newcommand{\tesisInstitucion}{Universidad de Buenos Aires \\ Facultad de Ciencias Exactas y Naturales}
\newcommand{\tesisInstitucionCorta}{FCEyN - UBA}
\newcommand{\tesisLugar}{Buenos Aires, 2026}
\newcommand{\tesisFecha}{Mayo de 2026}



\title{Plan de tesis de \tesisCarrera{\Huge \\ \tesisTitulo} }
\author{\tesisAutor\\{\small \tesisEmail\ - L.U. \tesisLU}}
\date{Direccion: \tesisDirector \\{\small \tesisDirectorEmail}~\vspace{0.2cm}\\ Codireccion: \tesisCodirector \\{\small \tesisCodirectorEmail}}


\begin{document}
 \maketitle

% \newpage
 
\section*{Introduccion y antecedentes de la tematica}

Describir aqui el area de investigacion, el problema general y los antecedentes mas relevantes. Reemplazar este texto por una presentacion breve del estado del arte y citar la bibliografia necesaria con \verb|\textcite{}| o \verb|\cite{}|.


\section*{Alcance, actividades y metodologia}

Describir el alcance concreto de la tesis, las preguntas que intenta responder y los limites del trabajo.

\subsection*{Actividades propuestas}

\begin{itemize}
    \item \textbf{Revision bibliografica}: relevar antecedentes y definir el marco teorico.
    \item \textbf{Preparacion de datos o materiales}: describir las fuentes, corpus, experimentos o insumos.
    \item \textbf{Desarrollo metodologico}: implementar, adaptar o analizar el metodo principal.
    \item \textbf{Evaluacion}: definir metricas, criterios de validacion o analisis de resultados.
    \item \textbf{Escritura}: redactar, revisar y preparar la entrega final.
\end{itemize}

\subsection*{Metodologia}
Describir las herramientas, datos, metodos y criterios de evaluacion que se usaran durante la tesis.

\section*{Plazos y factibilidad de realización}
El proyecto es factible para ser completado en el periodo previsto. El cronograma propuesto es:

\begin{table}[h]
\centering
\scriptsize 
\renewcommand{\arraystretch}{1.8}
\setlength{\tabcolsep}{2pt}
\begin{tabularx}{\textwidth}{|l|*{12}{>{\centering\arraybackslash}X|}}
\hline
\rowcolor{gray!20} 
\textbf{Actividad / Semana} & \textbf{1} & \textbf{2} & \textbf{3} & \textbf{4} & \textbf{5} & \textbf{6} & \textbf{7} & \textbf{8} & \textbf{9} & \textbf{10} & \textbf{11} & \textbf{12} \\ \hline
\textbf{Lunes (Inicio)} & 16/3 & 23/3 & 30/3 & 06/4 & 13/4 & 20/4 & 27/4 & 04/5 & 11/5 & 18/5 & 25/5 & 01/6 \\ \hline
Elaboración del plan & \cellcolor{gray!50}X & & & & & & & & & & & \\
\hline
Revisión bibliográfica y del estado del arte & & \cellcolor{gray!50}X &  & & & & & & & & & \\
\hline
Procesamiento de datos & & & \cellcolor{gray!50}X & \cellcolor{gray!50}X & \cellcolor{gray!50}X & & & & & & & \\
\hline
Implementación modelos & & & & & \cellcolor{gray!50}X & \cellcolor{gray!50}X & \cellcolor{gray!50}X & & & & & \\
\hline
Visualización y evaluación de los modelos & & & & & & & & \cellcolor{gray!50}X & \cellcolor{gray!50}X & & & \\ \hline
Entrenamiento y validación del modelo final & & & & & & & & & \cellcolor{gray!50}X & & & \\ \hline
Escritura del Manuscrito & & & & & & & & & \cellcolor{gray!50}X & \cellcolor{gray!50}X & \cellcolor{gray!50}X & \cellcolor{gray!50}X \\ \hline
\end{tabularx}
\caption{Cronograma de actividades.}
\end{table}

Describir aqui los recursos disponibles, dependencias externas y riesgos principales.

\newpage
\nocite{*}

\printbibliography

\end{document}
